\section{建设背景与目标}
\label{sec:background_goals}

\subsection{建设背景：企业大模型应用普及与数据治理需求}
近年来，生成式大语言模型（LLM）逐步融入企业日常经营。企业员工在软件研发、智能客服、文档翻译、合规审查及文案创作等场景中，广泛接入并使用多模态与代码大模型（例如 OpenAI、Kimi、DeepSeek 等）。

为了科学评估各业务部门的采纳情况、合理规划接口采购预算、优化模型负载调度，并对外部模型供应商的服务质量进行客观评估，企业数据治理团队需要定期统计并分析大模型的使用数据。典型的分析需求包括：
\begin{itemize}[leftmargin=2em]
  \item \textbf{部门渗透与使用广度：} 哪些业务部门已经开始使用大模型？日常活跃员工占比如何？
  \item \textbf{交互总量与成本消耗：} 各模型供应商的调用频次、Token 消耗总量及长周期增长趋势如何？
  \item \textbf{任务类型分布：} 员工主要利用模型解决哪些业务任务（如代码编写、翻译、文档总结）？
  \item \textbf{服务可用性与成功率：} 各供应商在复杂任务下的接口调用成功率与异常率如何？
\end{itemize}

\subsection{核心风险：传统统计方式的隐私与数据安全隐患}
在传统数据分析模式下，企业通常直接基于网关明细日志或聚合数据库生成报表。然而，针对大模型交互数据，直接发布明细数据或传统统计均值存在以下隐私与数据安全风险：
\begin{enumerate}[label=\textbf{\arabic*.}, leftmargin=2em]
  \item \textbf{高维行为特征可能导致个体被重新识别：} 大模型调用日志不仅包含员工标识，还记录了调用时间戳、任务类别、Token 长度、使用模型等丰富特征。即使隐去姓名工号进行脱敏，若结合外部公开信息（如考勤打卡时间、特定业务上线节点），依然存在通过高维特征碰撞反推特定员工行为轨迹的可能。
  \item \textbf{极值用户与长尾分布易扭曲大盘统计：} 企业内员工的使用习惯具有显著差异。少数编写自动化脚本的工程师单日可能产生数千甚至数万次调用，而大量普通员工单日仅有数次交互。如果不对单用户影响进行限制，极少数高频用户的行为会主导统计结果，且通过对比特定时段的数据变化，可能推断该员工是否参与了敏感项目。
  \item \textbf{外部数据共享中的业务机密保护：} 当企业向模型供应商反馈质量与成功率指标时，如果直接发布未经保护的细分数据，供应商可能间接推断出企业的核心业务规模与技术选型倾向。
\end{enumerate}

\subsection{建设目标与系统核心定位}
为兼顾“宏观数据可用”与“员工个体隐私”，\textbf{\mbox{Private LLM Analytics}} 系统被设计为企业大模型数据分析与合规发布的统一中间件。系统确立了以下核心建设目标：
\begin{itemize}[leftmargin=2em]
  \item \textbf{数学可证明的隐私保护：} 避免依赖脆弱的静态脱敏规则，引入差分隐私（Differential Privacy）作为底层机制。系统采用以单员工历史行为为边界的用户级保护（User-Level DP），在数学边界内有效限制单个员工记录对输出结果的可辨识影响。
  \item \textbf{高可用业务统计支持：} 在保障隐私预算的前提下，针对企业核心关注的调用量、Token、渗透率、成功率四项核心统计能力及扩展的高频任务等进行算法优化，使宏观统计误差保持在工程可接受的范围内，满足容量规划、财务对账与选型决策需求。
  \item \textbf{全生命周期预算闭环管理：} 建立企业级持久化隐私账本，跟踪控制跨时间、跨指标的累计隐私损耗；提供查询结果复用机制，降低重复查询带来的噪声平均风险。
  \item \textbf{多角色隔离与安全合规发布：} 建立供应商视角隔离与安全阈值抑制机制，支持安全导出供外部生态交互的合规分析报表。
\end{itemize}
